\section{H1 threshold decomposition and contribution closure}\label{supp:h1_bottleneck}

This descriptive analysis reads the existing exact action intervals, point bounds and paid baseline endpoints. It reads no additional physical response, evaluates no physical bound and changes no protocol. For branch threshold $T=B-\epsilon$, the main-text joint lower bound is at least $T$ exactly when the initial incumbent and both terms for every retained action are at least $T$. The proof is the elementary equivalence between a minimum exceeding a threshold and all of its arguments exceeding that threshold. The result describes sufficiency of the supplied relaxed certificate; it is not an impossibility result for the coupled GP or true engineering objective.

Both predeclared domains give71 nonbaseline contrasts and281 retained action terms. All71 contrasts have at least one pending term and one next-point term below threshold. There are71 failing pending terms and276 failing next-point terms. These are dependent descriptive counts; the two domains do not double the independent sample size. The median weakest next-point-bound/threshold ratio is0.1622233735904155. Raising only pending terms, or only next-point terms, cannot make the fixed-set certificate pass. Actual stronger physical information could change response domains or action sets, so no extrapolated impossibility claim is made.

The reader is \path{src/diagnose_h1_bound_bottlenecks.py}. Its output includes branch- and action-level rational thresholds, input SHA-256 hashes and a summary. Four synthetic checks cover the closed margin boundary, separate information bottlenecks, distinct blocking actions and the independent incumbent constraint. These checks contain no engineering optimization samples. The original31-test record remains separate.

The contribution judgment remains limited: numerical certificate production and its scoped computational improvement are supported; positive independent decision efficacy and full deployment savings are not. This additional diagnosis locates the physical-information bottleneck, but does not turn that negative outcome into an optimizer-performance claim or certify a journal novelty threshold.
